\documentclass{atomic}

\title{Products}
\newcommand{\breadcrumb}{Category Theory $>$ Universal Constructions}

\begin{document}
    \pagestyle{plain} \currentdoc{note} \noindent
    Let $\{ c_i \}_{i \in I}$ be a \ex{Collections}{family} of objects in a \ex{Categories}{category} $C$.
    \begin{definition}[Auxiliary Category] \label{auxiliary-category}
        The \textbf{auxiliary category} $D = D_\text{pr}$ is defined with
        \begin{itemize}
            \item objects being the pairs $\left( x, \{\alpha_i\}_{i \in I} \right)$, where $x \in C$ and $\alpha_i : x \to c_i$ for all $i \in I$,
            \item morphisms from an object $\left( x, \{\alpha_i\}_{i \in I} \right)$ to $\left( y, \{\beta_i\}_{i \in I} \right)$ being the morphisms $\gamma : x \to y$ in $C$ such that $\forall i \in I, \beta_i\gamma = \alpha_i$, i.e., such that the following diagram is \ex{CommutativeDiagrams}{commutative} for all $i$:
            $$\large\begin{tikzcd}
                x \arrow[dr, swap, "\alpha_i"] \arrow[rr, "\gamma"] & & y \arrow[dl, "\beta_i"] \\
                & c_i
            \end{tikzcd},$$
        \item and \ex{CompositeFunctions}{composition} being the same as in $C$.
        \end{itemize}
    \end{definition}

    A \ex{InitialAndTerminalObjects}{terminal object} $\left( p, \{ \pi_i \}_{i \in I} \right) \in D$ is called a \textit{product} of the \ex{Collections}{family} $\{ c_i \}_{i \in I}$. By \ex{InitialAndTerminalObjects}{definition}, for any $\left( x, \{\alpha_i\}_{i \in I} \right) \in D$, there exists a unique $\gamma : x \to p$ in $C$ such that $\forall i, \pi_i\gamma = \alpha_i$, i.e., such that the following diagram is \ex{CommutativeDiagrams}{commutative} for all $i$:
    $$\large\begin{tikzcd}
        x \arrow[dr, swap, "\alpha_i"] \arrow[rr, "\exists!\gamma"] & & p \arrow[dl, "\pi_i"] \\
        & c_i
    \end{tikzcd}.$$

    \begin{remark*}
        By \ex[uniqueness]{InitialAndTerminalObjects}{the uniqueness of terminal objects up to} \ex{Isomorphisms}{isomorphism}, \textit{products} (if they exist) are unique up to \ex{Isomorphisms}{isomorphism} in $D$, which amounts to \ex{Isomorphisms}{isomorphism} in $C$ that preserves the \ex{CommutativeDiagrams}{commutativity}.
    \end{remark*}

    \begin{definition}[``Product Exists'']
        We say that a ``\textbf{product exists}'' in $C$ iff a \textit{product} exists for every \ex{Collections}{family} of objects in $C$.
    \end{definition}

    \section{Binary Product} \label{binary-product}
    $$\large\begin{tikzcd}[column sep={2em}, row sep={2em}]
        & x \arrow[d, dashed, "\exists!f"] \\
        & \arrow[dl, "\pi_1"] a \times b \arrow[dr, swap, "\pi_2"] \\[-2em]
        a && b
        \arrow[from=1-2, to=3-1, bend right, "f_1"']
        \arrow[from=1-2, to=3-3, bend left, "f_2"]
    \end{tikzcd}$$
    When $\{ c_i \}_{i \in I} = \{ a,b \}$, and $p$ is denoted $a \times b$, you get the above. Stated without the language of the \hyperref[auxiliary-category]{auxiliary category}:

    \begin{quote}
        ``Objects $a,b \in C$ have a \textit{product} $a \times b$ and $\pi_1,\pi_2$ iff for every $x \in C$, there exists a unique morphism $f : x \to (a \times b)$ such that the above diagram is \ex{CommutativeDiagrams}{commutative} for all $f_1$ and $f_2$.''
    \end{quote}

\end{document}
